\section*{Hoofdstuk \ref{ch:b2:charges-currents-conduction} --- Ladingen,
stromen en geleiding}

\begin{solution}{exo:b2:charges-currents-conduction:1}
$j = 16/2.5 \times 10^{-6} = \qty{6.4e6}{A/m^2}$;
$v = j/ne = \qty{0.47}{mm/s}$; \qty{20}{m} in $4.3 \times 10^4$ s, twaalf
uur. $E = j/\gamma = \qty{0.11}{V/m}$; $U = \qty{2.1}{V}$;
$P = UI = \qty{34}{W}$.
\end{solution}

\begin{solution}{exo:b2:charges-currents-conduction:2}
Bliksem $3 \times 10^4/\pi10^{-4} = \qty{1e8}{A/m^2}$; bundel
$\qty{4e3}{A/m^2}$; axon $10^{-9}/1.3 \times 10^{-11} = \qty{80}{A/m^2}$;
ringstroom \qty{1e-8}{A/m^2}. Alleen de bliksem overtreft de $6 \times 10^6$
van de draad.
\end{solution}

\begin{solution}{exo:b2:charges-currents-conduction:3}
$R = \rho_eL/S$: $0.68$, $1.1$, $4.0$, \qty{44}{\ohm}. Aluminium:
$2.5 \times 2.8/1.7 =
\qty{4.1}{mm^2}$, massaverhouding
$(4.1 \times 2700)/(2.5 \times 8900) = 0.50$ --- de helft van de massa.
Verwarming: $R = 230^2/1000 = \qty{53}{\ohm}$, $S = \qty{0.196}{mm^2}$,
$L = RS/\rho_e = \qty{9.4}{m}$.
\end{solution}

\begin{solution}{exo:b2:charges-currents-conduction:4}
$\tau = m\gamma/ne^2$: koper \qty{2.5e-14}{s},
$\mu = e\tau/m = \qty{4.4e-3}{m^2/(V.s)}$, $\ell = \qty{40}{nm}$; zilver
\qty{3.8e-14}{s}, \qty{6.7e-3}{m^2/(V.s)}, \qty{61}{nm}. Silicium:
$\gamma = ne\mu = 10^{22} \times 1.6 \times 10^{-19} \times 0.14 = \qty{220}{S/m}$;
$v =
\mu E = \qty{140}{m/s}$.
\end{solution}

\begin{solution}{exo:b2:charges-currents-conduction:5}
(a) Er komt stroom binnen en er gaat er geen uit: $\dd Q/\dd t = I$. (b) Er
hoopt zich lading op: $\partial_t
\rho \ne 0$. (c)
$\vect j = (I/4\pi r^2)\vect e_r$:
$\operatorname{div}\vect j = \frac1{r^2}
\partial_r(I/4\pi) = 0$. (d) Alleen
bij de elektrode, waar de stroom vanuit de draad wordt ingespoten --- een
bron van de flux.
\end{solution}

\begin{solution}{exo:b2:charges-currents-conduction:6}
(a) $j = I/2\pi rL$, $E = j/\gamma$, radiaal. (b)
$U = \int_a^bE\,\dd r = I\ln(b/a)/
2\pi\gamma L$:
$R = \ln3.5/(2\pi \times 10^{-13} \times 10^3) = \qty{2e9}{\ohm}$ per
kilometer. (c) \qty{0.5}{\micro A}, \qty{0.5}{mW}. (d)
$RC = \varepsilon_0\varepsilon_r/\gamma = \qty{200}{s}$ voor
$\varepsilon_r = 2.3$: de geladen en geïsoleerde kabel ontlaadt zich door
haar eigen isolatie met de relaxatietijd van die isolatie.
\end{solution}

\begin{solution}{exo:b2:charges-currents-conduction:7}
(a) $\omega\tau$: $8 \times 10^{-12}$, $1.6 \times 10^{-3}$, $4.7$, $79$.
(b) $|\underline\gamma|/\gamma_0 =
1/\sqrt{1 + \omega^2\tau^2}$: $1$, $1$,
$0.21$, $0.013$; fase $-\arctan\omega\tau$: $0$, $-\qty{0.09}{\degree}$,
$-\qty{78}{\degree}$, $-\qty{89}{\degree}$. (c) $\omega\tau \le 0.14$:
$f \le
\qty{0.9}{THz}$. (d) $\underline\gamma \approx ne^2/\iu m\omega$:
$\underline{\vect j}$ loopt \qty{90}{\degree} achter op
$\underline{\vect E}$, $\langle\vect j\cdot\vect E\rangle = 0$: geen
warmteontwikkeling --- de elektronen trillen vrij, zoals in een plasma.
\end{solution}

\begin{solution}{exo:b2:charges-currents-conduction:8}
(a) $S = \qty{3.1e-8}{m^2}$, $j = \qty{3.2e8}{A/m^2}$,
$p = j^2/\gamma = \qty{1.7e9}{W/m^3}$. (b) Energie om te smelten per
volume-eenheid
$\rho(c\Delta T + L_f) = 8900(4.1 \times 10^5 + 2.05 \times
10^5) = \qty{5.5e9}{J/m^3}$:
\qty{3.2}{s}. (c) Negen keer sneller: \qty{0.36}{s}. (d) De nominale stroom
is die waarvan de lucht de warmte van Joule net onder het smeltpunt kan
afvoeren; het zand blust de vlamboog die anders na het smelten van de draad
over de spleet zou blijven geleiden.
\end{solution}

\begin{solution}{exo:b2:charges-currents-conduction:9}
(a)
$\partial_t\rho = -\operatorname{div}\vect j = -\gamma\operatorname{div}\vect E = -\gamma\rho/
\varepsilon_0$.
(b) $\varepsilon_0/\gamma$: \qty{1.5e-19}{s}, \qty{1.8e-12}{s},
\qty{1.8e-10}{s}, \qty{9}{s}, $9 \times 10^4$ s. (c)
$f \lesssim 1/2\pi\tau_r$: $10^{18}$ Hz (het \omterm{prop:b2:charges-currents-conduction:drude}{model van Drude} faalt lang
daarvoor al), \qty{90}{GHz}, \qty{0.9}{GHz}, \qty{0.02}{Hz}, \qty{2e-6}{Hz}.
(d) Een glazen staaf zou haar lading in een minuut door haar volume
verliezen; echt droog glas zit eerder bij $10^{-14}$ S/m (uren), en meestal
is het de vochtigheid aan het oppervlak die haar ontlaadt --- in orde van
grootte klopt het.
\end{solution}

\begin{solution}{exo:b2:charges-currents-conduction:10}
(a) $\vect j = (I/2\pi r^2)\vect e_r$, $E = I/2\pi\gamma r^2$,
$V = I/2\pi\gamma r$. (b) $R = V(a)
/I = 1/2\pi\gamma a = \qty{32}{\ohm}$.
(c) $V_{\text{staaf}} = \qty{640}{kV}$; $V(10) = \qty{32}{kV}$,
$V(10.7) = \qty{30}{kV}$: \qty{2.1}{kV} tussen de voeten. (d) Voor- en
achterpoten liggen \qty{1.5}{m} uit elkaar langs de straal:
$4$--\qty{5}{kV}, en de stroomweg loopt door het hart.
\end{solution}

\begin{solution}{exo:b2:charges-currents-conduction:11}
(a) $E_H = jB/nq$ met $j = I/wt$: $V_H = E_Hw = IB/nqt$. (b)
$1/(8.5 \times 10^{28}
\times 1.6 \times 10^{-19} \times 10^{-4}) = \qty{0.7}{\micro V}$.
(c) \qty{6.3}{mV}: $V_H \propto 1/n$, en halfgeleiders hebben $10^6$ keer
minder dragers. (d)
$B = V_Hnqt/I =
10^{-5} \times 10^{22} \times 1.6 \times 10^{-19} \times 2 \times 10^{-5}/10^{-3} = \qty{0.32}{T}$;
het teken van $V_H$ geeft het teken van de dragers (bij bekende $\vect B$),
of de richting van $\vect B$ (bij een bekende sensor).
\end{solution}

\begin{solution}{exo:b2:charges-currents-conduction:12}
(a) $\vect j = ne\mu_+\vect E + (-ne)(-\mu_-\vect E)$: beide soorten driften
in tegengestelde richtingen en voeren stroom in dezelfde richting. (b)
$n =
0.5 \times 10^3 \times 6.02 \times 10^{23} = \qty{3e26}{m^{-3}}$:
$\gamma = 3 \times 10^{26} \times 1.6 \times 10^{-19}
\times 1.31 \times 10^{-7} = \qty{6.3}{S/m}$
(bij deze concentratie hinderen de ionen elkaar: \qty{5}{S/m} gemeten). (c)
\qty{0.52}{\micro m/s} en \qty{0.79}{\micro m/s}: Cl$^-$ voert $60\%$. (d)
$R = L/\gamma S = \qty{200}{\ohm}$, $U = \qty{200}{V}$, $P =
\qty{200}{W}$
in \qty{10}{g} water: \qty{10}{K} in \qty{2}{s} --- vandaar de koelplaten
van elektroforesebakken.
\end{solution}

\begin{solution}{pb:b2:charges-currents-conduction:1}
\textbf{1.} $n = \rho N_A/M = \qty{8.4e28}{m^{-3}}$.

\textbf{2.} $j = \qty{6.4e6}{A/m^2}$, $v = j/ne = \qty{0.47}{mm/s}$;
\qty{35}{min} per meter.

\textbf{3.} $\tau = m\gamma/ne^2 = \qty{2.5e-14}{s}$,
$\mu = \qty{4.4e-3}{m^2/(V.s)}$, $\ell =
\qty{40}{nm}$, $150$
atoomafstanden: de elektronen botsen tegen defecten en thermische
trillingen, niet tegen de atomen.

\textbf{4.} $\rho_e \propto T$ geeft
$\alpha = 1/293 = \qty{3.4e-3}{K^{-1}}$, dicht bij de gemeten
$4 \times 10^{-3}$.

\textbf{5.} $E = \qty{0.11}{V/m}$, $U = \qty{0.11}{V}$;
$\tfrac12mv^2 = \qty{1e-37}{J}$ tegen $k_BT = \qty{4e-21}{J}$: het
elektronengas merkt niets van de stroom.

\textbf{6.} $I/e = \qty{1e20}{}$ elektronen per seconde.

\textbf{7.} $Q = neSL = \qty{3.4e4}{C}$; $QE = \qty{3.6}{kN}$, tegen een
gewicht van \qty{0.2}{N}. Het veld trekt met de tegengestelde kracht aan het
positieve rooster; de elektronen geven hun impuls door botsingen aan het
rooster af: de draad als geheel voelt niets.

\textbf{8.} $p = j^2/\gamma = \qty{6.8e5}{W/m^3}$, \qty{1.7}{W/m}.

\textbf{9.} Oppervlak $\pi d = \qty{5.6e-3}{m^2}$ per meter:
$\Delta T = 1.7/(10 \times 5.6 \times
10^{-3}) = \qty{30}{K}$; de
\omterm{prop:b2:charges-currents-conduction:drude}{soortelijke weerstand} stijgt $12\%$ en het vermogen tot \qty{1.9}{W}:
$\Delta T \approx \qty{35}{K}$.

\textbf{10.}
$\Delta T = 1.7(1 + 0.004\Delta T)/(2 \times 5.6 \times 10^{-3})$:
$\Delta T(1 - 0.6) = 150$, $\Delta T \approx \qty{400}{K}$ --- dicht bij
thermisch op hol slaan: ingegraven kabels moeten lager worden belast.

\textbf{11.} $p/\rho c = \qty{0.2}{K/s}$; bij \qty{160}{A} \qty{20}{K/s}:
smelten in ongeveer een minuut (minder naarmate $\rho_e$ stijgt); de
automaat schakelt binnen milliseconden uit.

\textbf{12.} $\omega\tau = 8 \times 10^{-12}$: de wet van Ohm blijft onaangetast.

\textbf{13.} $\varepsilon_0/\gamma = \qty{1.5e-19}{s}$: geen lading
binnenin; het veld in de draad wordt gemaakt door ladingen op haar
\emph{oppervlak}, waarvan de dichtheid langs de draad verandert.

\textbf{14.} $Q = I_0\tau_c = \qty{16}{mC}$, $V = Q/C = \qty{160}{V}$.

\textbf{15.} Aan de draadzijde loopt $I$, aan de spleetzijde geen
geleidingsstroom: de lading groeit, en met haar het elektrische veld in de
spleet.

\textbf{16.} $I = 16\eu^{-0.5} = \qty{9.7}{A}$; $E = Q/\varepsilon_0A$,
$\dd E/\dd t = I/\varepsilon_0A
= 9.7/(8.85 \times 10^{-12} \times 10^{-4}) = \qty{1.1e16}{V.m^{-1}.s^{-1}}$.

\textbf{17.} $\varepsilon_0A\,\dd E/\dd t = I$: de verplaatsingsstroom.

\textbf{18.}
$j = \gamma E = 10^{-14} \times 2.3 \times 10^5 = \qty{2e-9}{A/m^2}$,
$10^{-14}$ A door de doorsnede: niets; doorslag bij \qty{3}{kV} over de
millimeter.

\textbf{19.} $R = 1000/(6 \times 10^7 \times 10^{-5}) = \qty{1.7}{\ohm}$;
$\Delta U = \qty{67}{V}$; $P =
\qty{2.7}{kW}$, $29\%$ van de geleverde
\qty{9.2}{kW}.

\textbf{20.} $I = \qty{0.46}{A}$, $\Delta U = \qty{0.8}{V}$, \qty{0.35}{W}:
transport bij hoge spanning deelt het verlies door $(U_2/U_1)^2$.

\textbf{21.} $10 \times 6/3.6 = \qty{17}{mm^2}$; \qty{89}{kg/km} koper tegen
\qty{45}{kg/km} aluminium: lichter en goedkoper op masten; in de muur past
de kleinere doorsnede van koper beter in de buizen en klemmen.

\textbf{22.}
$V_H = IB/nqt = 5 \times 10^{-3} \times 0.1/(10^{22} \times 1.6 \times 10^{-19} \times 10^{-4}) =
\qty{3.1}{mV}$.

\textbf{23.} $+24\%$: \qty{2.1}{\ohm}, \qty{3.3}{kW} verlies.

\textbf{24.}
$RI^2t = 1.7 \times 4 \times 10^8 \times 10^{-4} = \qty{67}{kJ}$ in
\qty{89}{kg}: \qty{2}{K}. $j = \qty{2e9}{A/m^2}$, driehonderd keer de
huishoudelijke waarde.

\textbf{25.} $n$ (de materie) en $\tau$ (de botsingen) geven
$\gamma = ne^2\tau/m$; $\gamma$ verbindt $j$ en $E$ (Ohm); $j$ integreert
tot $I$ en $E$ tot $U$ (de schakeling); $j^2/\gamma$ is de warmte.
\end{solution}
